\documentclass[11pt]{article}
\usepackage[T1]{fontenc}
\usepackage{lmodern}
\usepackage[utf8]{inputenc}
\usepackage[a4paper,margin=25mm]{geometry}
\usepackage{amsmath,amssymb,amsthm,mathtools}
\usepackage{microtype}
\usepackage{enumitem}
\usepackage{hyperref}
\hypersetup{colorlinks=true,linkcolor=blue,citecolor=blue,urlcolor=blue,
  pdftitle={First correction to the critical mean degree in extreme long-range percolation},
  pdfauthor={Anonymous},pdfsubject={Unrefereed research candidate}}
\newtheorem{theorem}{Theorem}[section]
\newtheorem{proposition}[theorem]{Proposition}
\newtheorem{lemma}[theorem]{Lemma}
\newtheorem{remark}[theorem]{Remark}
\theoremstyle{definition}
\newtheorem{definition}[theorem]{Definition}
\newcommand{\Z}{\mathbb Z}
\newcommand{\Pp}{\mathbb P_p}
\newcommand{\Ep}{\mathbb E_p}
\newcommand{\dd}{\,\mathrm d}
\newcommand{\ind}{\mathbf 1}
\newcommand{\del}{\mathrm{del}}
\newcommand{\NB}{\mathrm{NB}}
\newcommand{\through}[3]{\{#1\xleftrightarrow{\,#3\,}#2\}}
\newcommand{\intk}{\int_{-\pi}^{\pi}}
\newcommand{\dm}{\frac{\mathrm dk}{2\pi}}
\DeclareMathOperator{\supp}{supp}
\title{First correction to the critical mean degree\protect\\
in extreme long-range percolation}
\author{Anonymous}
\date{11 October 2026}

\begin{document}
\maketitle

\begin{abstract}
Consider independent undirected bond percolation on the integer lattice, with
connection probability $C|x-y|^{-1-\sigma}$, where $0<\sigma<1$. We determine
the first correction to the critical expected degree as $\sigma$ tends to zero.
Writing $p_c=2\zeta(1+\sigma)C_c(\sigma)$, we prove
\[
 p_c=1+\frac{\zeta(2+2\sigma)}{2\zeta(1+\sigma)^2}+O(\sigma^3)
     =1+\frac{\pi^2}{12}\sigma^2+O(\sigma^3).
\]
The error is uniform for all sufficiently small positive real $\sigma$, with
dimension fixed and no additional spread parameter. The lower bound follows
from an exact nonbacktracking path estimate. For the upper bound, elementary
Fourier estimates verify the general small-triangle hypotheses for the lace
expansion. We then retain a deleted-bond restriction in the first lace
coefficient and distinguish a closed diagram of order $\sigma^2$ from its
spatial supremum of order $\sigma$. The positive remainder after the first odd
truncation converts a uniform subcritical coefficient bound into the threshold
bound. This avoids a critical limit of lace coefficients and an estimate of
the entire higher-order lace tail in the additional correction calculation;
the imported infrared theorem itself still relies on lace-expansion estimates.
We extract a two-scale kernel criterion and apply it uniformly to arbitrary
nonnegative summable additions and compact mixtures of nearby tail exponents.
For an added nearest-neighbor weight $\lambda$, the coefficient becomes
$\pi^2/12+\lambda+\lambda^2/2$; for exponent mixtures it is multiplied by
the square of the mean exponent ratio.
\end{abstract}

\noindent\textbf{Keywords:} long-range percolation; critical probability;
lace expansion; nonbacktracking paths; small-exponent asymptotics.

\section{Model and result}

For $0<\sigma<1$ and $0\le C\le1$, declare each unordered pair
$\{x,y\}$ of distinct integers open independently with probability
$C|x-y|^{-1-\sigma}$. Let $C_c(\sigma)$ be the infimum of the parameters for
which the origin belongs to an infinite open cluster with positive probability.
At $C=1$ all nearest-neighbor bonds are open, so this infimum is taken over a
nonempty set. Put
\begin{equation}\label{eq:kernel}
 Z_\sigma=\zeta(1+\sigma),\qquad
 D_\sigma(0)=0,\qquad
 D_\sigma(x)=\frac{|x|^{-1-\sigma}}{2Z_\sigma}\quad(x\ne0).
\end{equation}
Thus $D_\sigma$ is an even probability distribution. The normalized parameter
$p=2Z_\sigma C$ is the expected degree, and the bond probabilities become
$pD_\sigma(y-x)$. Its admissible interval is
$[0,\|D_\sigma\|_\infty^{-1}]=[0,2Z_\sigma]$. We write
$p_c(\sigma)=2Z_\sigma C_c(\sigma)$.

\begin{theorem}\label{thm:main}
There are constants $K<\infty$ and $\sigma_0>0$ such that, for every real
$0<\sigma<\sigma_0$, the model defined above satisfies
\begin{equation}\label{eq:main}
 \left|p_c(\sigma)-1-S_2(\sigma)\right|\le K\sigma^3,
 \qquad
 S_2(\sigma)=\frac{\zeta(2+2\sigma)}{2\zeta(1+\sigma)^2}.
\end{equation}
In particular,
\begin{equation}\label{eq:limit}
 \lim_{\sigma\downarrow0}
 \frac{2\zeta(1+\sigma)C_c(\sigma)-1}{\sigma^2}
 =\frac{\pi^2}{12}.
\end{equation}
\end{theorem}

In the original bond-amplitude parameter the same result reads
\begin{equation}\label{eq:amplitude}
 C_c=\frac{1}{2Z_\sigma}
       +\frac{\zeta(2+2\sigma)}{4Z_\sigma^3}+O(\sigma^4),
 \qquad
 C_c-\frac{1}{2Z_\sigma}
       =\frac{\pi^2}{24}\sigma^3+O(\sigma^4).
\end{equation}
Indeed, divide \eqref{eq:main} by $2Z_\sigma$ and use
$Z_\sigma^{-1}=\sigma+O(\sigma^2)$. This is an asymptotic improvement
over the branching bound, not an evaluated error guarantee at a specified
positive exponent.

The leading limit $p_c\to1$ is raised for this undirected model in
\cite[Section~7]{Amit2026}. That paper develops explicit upper bounds using
monotone paths, which also apply to the oriented problem. The correction in
\eqref{eq:main} concerns the undirected model and its exact linear Bernoulli
bond probabilities.

The term $D^{*2}(0)$ is familiar from critical-point expansions. In a
spread-out regime above the upper critical dimension, van der Hofstad and
Sakai \cite[Theorem~1.1, (1.14)]{HS2005} obtain
\[
 p_c=1+D^{*2}(0)+\frac12\sum_{n\ge3}(n+1)D^{*n}(0)+O(\beta^2),
 \qquad \beta=L^{-d}.
\]
That theorem fixes the profile and takes the spread parameter $L$ large;
its percolation statement assumes $d>6$. It does not directly cover
\eqref{eq:kernel} with dimension one and $\sigma\downarrow0$. Furthermore,
the general lace-expansion parameter here must dominate
$\|D_\sigma\|_\infty=\Theta(\sigma)$, so an error of its square would be
too large to determine \eqref{eq:limit}.

We use the kernel-general infrared theorem of Heydenreich, van der Hofstad
and Sakai \cite{HHS2008}, and the first-coefficient events of
Borgs, Chayes, van der Hofstad, Slade and Spencer \cite{BCHSS2005}.
The additional estimates below keep the dependence on $\sigma$ explicit.
The two restrictions that matter are the original first attachment being
nonzero, and a terminal vertex lying outside a deleted-bond cluster. They
leave a single main term in which deletion of the first bond removes one
of two apparent two-step contributions. Only a one-sided bound on this
coefficient is required.

Throughout, $f*g$ denotes convolution on $\Z$, $f^{*m}$ a convolution
power, and $\delta(x)=\ind_{\{x=0\}}$. Fourier transforms are
$\widehat f(k)=\sum_x f(x)e^{ikx}$, with normalized measure $\mathrm dk/(2\pi)$
on $[-\pi,\pi]$. All constants denoted by $c$ or $K$ are finite and may
increase from line to line. Unless additional indices are shown, they are
independent of $\sigma$ and of $p\in[1,p_c(\sigma))$ on one sufficiently
small fixed interval $0<\sigma<\sigma_0$. No claim of uniformity in growing
convolution orders is made.

\section{Uniform estimates for the power-law kernel}\label{sec:kernel}

Write $s=\sigma$, $Z=Z_s$, $D=D_s$, and $d(k)=\widehat D_s(k)$ within proofs.
Integral comparison gives
\begin{equation}\label{eq:zeta-bounds}
 s^{-1}\le Z\le1+s^{-1}.
\end{equation}
In particular $\|D\|_\infty\le s/2$. The following estimates also control
the very small Fourier frequencies, where pointwise convergence of $d$ to
zero would not suffice.

\begin{lemma}\label{lem:fourier}
Let $L_0=\pi+1+\log2$. For $0<s\le1/2$ and $0<|k|\le\pi$,
\begin{align}
 |d(k)|&\le \min\{1,s[L_0+\log_+(1/|k|)]\},\label{eq:log-bound}\\
 1-d(k)&\ge\frac13\left(\frac{|k|}{2\pi}\right)^s.\label{eq:gap}
\end{align}
For fixed integers $r\ge1$, $q\ge1$, there is a constant $K_{r,q}$ such
that
\begin{equation}\label{eq:J}
 J_{r,q}(s):=\intk\frac{|d(k)|^r}{[1-d(k)]^q}\dm
 \le K_{r,q}s^r
 \quad\left(0<s\le\min\{1/2,1/(2q)\}\right).
\end{equation}
The corresponding estimate with $q=0$ holds for $0<s\le1/2$.
\end{lemma}

\begin{proof}
By symmetry it is enough to consider $0<k\le\pi$. Partial sums of
$e^{ink}$ over an interval have modulus at most $1/\sin(k/2)$. Abel
summation against the decreasing sequence $n^{-1-s}$ therefore gives
\begin{equation}\label{eq:abel}
 \left|\sum_{n=N}^\infty n^{-1-s}e^{ink}\right|
 \le\frac{N^{-1-s}}{\sin(k/2)}
 \le\frac{\pi}{k}N^{-1-s}\qquad(N\ge1).
\end{equation}
For $k\le1$, split the series defining $Zd(k)$ at $N=\lceil1/k\rceil$.
The head is at most $1+\log N\le1+\log2+\log(1/k)$ in absolute value,
and \eqref{eq:abel} bounds the tail by $\pi$. For $k\ge1$, use
\eqref{eq:abel} with $N=1$. Division by $Z\ge1/s$ gives the second
bound in \eqref{eq:log-bound}; the bound by one follows because $d$ is
a characteristic function.

For the spectral gap, take $N=\lceil\pi/k\rceil$, so that
$\pi/k\le N\le2\pi/k$. Positivity of $1-\cos(nk)$ and the tail integral
give
\begin{align*}
 Z[1-d(k)]
 &\ge\sum_{n=N}^\infty n^{-1-s}[1-\cos(nk)]\\
 &\ge \frac{N^{-s}}s-\frac{\pi}{k}N^{-1-s}
 \ge N^{-s}(s^{-1}-1).
\end{align*}
Since $(1-s)/(sZ)\ge(1-s)/(1+s)\ge1/3$, this proves
\eqref{eq:gap}. Retaining the entire tail is what avoids a loss of a
factor $s$ in this estimate.

Combining the two bounds shows that
\[
 J_{r,q}(s)\le
 \frac{3^q(2\pi)^{qs}}{\pi}s^r
 \int_0^\pi[L_0+\log_+(1/k)]^r k^{-qs}\dd k.
\]
On $(0,1)$, the substitution $t=-\log k$ bounds the last integral by
$\int_0^\infty(L_0+t)^r e^{-t/2}\dd t$ when $qs\le1/2$. On $[1,\pi]$
it is at most $\pi L_0^r$. Also $(2\pi)^{qs}\le(2\pi)^{1/2}$.
These bounds provide an explicit finite choice of $K_{r,q}$.
For $q=0$ the same calculation omits the gap factor.
\end{proof}

We shall need the following immediate consequences:
\begin{align}
 S_2=D^{*2}(0)=\intk d(k)^2\dm
 &=\frac{\zeta(2+2s)}{2Z^2}
   =\frac{\pi^2}{12}s^2+O(s^3),\label{eq:S2}\\
 I_{2,1}:=\intk\frac{d(k)^2}{1-d(k)}\dm
 &=S_2+O(s^3),\label{eq:I21}\\
 J_{2,3}&=O(s^2),\qquad J_{3,2}=O(s^3).\label{eq:needed-J}
\end{align}
For \eqref{eq:S2}, \eqref{eq:zeta-bounds} implies $Z^{-1}=s+O(s^2)$,
and the convergent series $\sum_{n\ge1}(\log n)n^{-2}$ gives
$\zeta(2+2s)=\zeta(2)+O(s)$. For \eqref{eq:I21}, use
\[
 I_{2,1}-S_2=\intk\frac{d(k)^3}{1-d(k)}\dm,
 \qquad |I_{2,1}-S_2|\le2J_{3,2},
\]
since $1-d(k)\le2$. None of these estimates assumes $d(k)\ge0$.

\section{A nonbacktracking lower bound}\label{sec:lower}

For this section define the pointwise moments
\[
 S_m=\sum_{x\ne0}D(x)^m
 =2^{1-m}\frac{\zeta(m(1+s))}{Z^m}\qquad(m\ge2).
\]
Only $S_2$ is identified here with a convolution return probability;
in particular, $S_3=\sum_xD(x)^3$ is not $D^{*3}(0)$.

\begin{proposition}\label{prop:NB}
For sufficiently small $s>0$, the equation
\begin{equation}\label{eq:NB-root}
 F_s(p):=\sum_{x\ne0}\frac{pD(x)}{1+pD(x)}=1
\end{equation}
has a unique solution $p_\NB\in(1,2)$. Moreover,
\begin{equation}\label{eq:NB-result}
 p_c\ge p_\NB=1+S_2+O(S_3+S_2^2)
              =1+S_2+O(s^3).
\end{equation}
\end{proposition}

\begin{proof}
Fix $p>0$ with $F:=F_s(p)<1$, and put $q_x=pD(x)$ and
$f(x)=(1+q_x)^{-1}$ for $x\ne0$. On bounded nonnegative functions of a
previous displacement define the positive operator
\[
 (Bh)(x)=\sum_{y\ne0,\,y\ne-x}q_yh(y).
\]
It is bounded in the supremum norm by $p$. Symmetry gives
\[
 Bf(x)=F-\frac{q_x}{1+q_x}=f(x)-(1-F)\le Ff(x),
\]
where the last inequality uses $f(x)\le1$. Writing
$q_{\max}=\sup_xq_x$, we have $1\le(1+q_{\max})f$, so positivity and
induction imply
\[
 B^{n-1}1\le(1+q_{\max})F^{n-1}f.
\]
The total product weight $W_n$ of nonbacktracking walks of length $n$
starting at zero consequently satisfies
\begin{equation}\label{eq:walk-weight}
 W_n=\sum_xq_x(B^{n-1}1)(x)
 \le(1+q_{\max})F^n\qquad(n\ge1).
\end{equation}
All these sums are nonnegative and finite; for example $W_n\le p^n$.

Every self-avoiding path is nonbacktracking, and its bond-open probability
is the product of its bond probabilities because its unordered bonds are
distinct. A union bound therefore bounds the probability of an open
self-avoiding path of length $n$ from zero by $W_n$. We do not assign
product probabilities to the nonbacktracking walks that repeat bonds;
their weights only enlarge the numerical upper bound. Each vertex has
finite open degree almost surely, since its expected degree is $p<\infty$.
By countability the random graph is locally finite almost surely. An infinite
component then has self-avoiding paths of every finite length from each of
its vertices. Equation~\eqref{eq:walk-weight} proves absence of an infinite
cluster whenever $F_s(p)<1$. The case $p=0$ is immediate.

The function $F_s$ is continuous and strictly increasing, by domination
by $pD(x)$ on bounded intervals. The exact expansion
\begin{equation}\label{eq:NB-expansion}
 F_s(p)=p-p^2S_2+p^3T_3(p),\qquad
 T_3(p)=\sum_x\frac{D(x)^3}{1+pD(x)},\quad0\le T_3(p)\le S_3,
\end{equation}
shows $F_s(1)<1$. If $a=S_2\le1/8$, then
\[
 F_s(1+2a)\ge1+2a-(1+2a)^2a
 =1+a(1-4a-4a^2)>1.
\]
There is thus exactly one root in $(1,1+2a)$, and $p_c\ge p_\NB$.
The interval lies below $2<2Z$, so all parameters used are admissible.
For instance $0<s\le1/4$ ensures $S_2\le\zeta(2)s^2/2<1/8$.

Set $x=p_\NB-1\in(0,2a)$. Equation~\eqref{eq:NB-expansion} yields
\[
 x=p_\NB^2a-p_\NB^3T_3(p_\NB),\qquad
 |x-a|\le(4a+4a^2)a+8S_3.
\]
Since $S_3\le\zeta(3)s^3/4$, the asserted uniform expansion follows.
\end{proof}

\section{The infrared input and the required diagram estimates}
\label{sec:diagrams}

For fixed $s$, write
$\tau_p(x)=\Pp(0\leftrightarrow x)$ and
$\chi(p)=\sum_x\tau_p(x)$. We use the standard equality of the
susceptibility threshold and the percolation threshold, in the form stated
for this long-range model in \cite[Section~1.2.5]{HHS2008}. Consequently
$\chi(p)<\infty$ for every $p<p_c$. This sharpness input is distinct
from a definition of $p_c$ by susceptibility.

For percolation, Assumption~2.1 of \cite{HHS2008} requires a symmetric
probability kernel with $D(0)=0$, together with
\begin{equation}\label{eq:HHS-assumptions}
 \|D\|_\infty\le\beta,\qquad
 \intk\frac{d(k)^2}{[1-d(k)]^3}\dm\le\beta,
\end{equation}
for sufficiently small $\beta$. The kernel \eqref{eq:kernel} has full
support off zero, in particular it contains the nearest neighbors.
Lemma~\ref{lem:fourier} verifies \eqref{eq:HHS-assumptions} with
$\beta=B_0s$ for a fixed $B_0$ and all sufficiently small $s$.
Here $B_0$ is independent of $s$; it is chosen before reducing the interval
of $s$. Theorem~2.3, equation~(1.35), and the threshold bound~(2.6) of
\cite{HHS2008} then give constants $K_0,K_1$ such that
\begin{align}
 1\le p_c&\le1+K_0s,\label{eq:coarse-pc}\\
 0\le\widehat\tau_p(k)&\le
 \frac{1+K_1s}{\chi(p)^{-1}+p[1-d(k)]},\qquad p<p_c.
 \label{eq:IR}
\end{align}
The constants in this use of the theorem are fixed by its smallness
threshold and the dimension $d=1$, not by an additional spread parameter.
The parameter $s=3$ in its percolation assumption is a diagram exponent;
it is unrelated to our variable $s=\sigma$.

For completeness, the Fourier nonnegativity in \eqref{eq:IR} is a
property of percolation connections. For every finitely supported complex
sequence $(a_x)$,
\[
 \sum_{x,y}a_x\overline{a_y}\tau_p(y-x)
 =\Ep\sum_{\mathcal C}
       \left|\sum_{x\in\mathcal C}a_x\right|^2\ge0,
\]
where the sum is over the open clusters. Thus $\tau_p$ is positive
definite; since it is in $\ell^1$ subcritically, its continuous Fourier
transform is nonnegative. This does not assert nonnegativity of $d(k)$.
For the rest of the proof we restrict to $1\le p<p_c$, where
\begin{equation}\label{eq:IR-simplified}
 \widehat\tau_p(k)\le\frac{b_s}{p[1-d(k)]},\qquad b_s=1+K_1s,
 \quad k\ne0.
\end{equation}
The single point $k=0$ has no effect on the integrals used below.

Suppress $p$ from $\tau_p$ and put $q=pD$. Define
\begin{equation}\label{eq:diagram-defs}
 T(x)=(q*\tau^{*3})(x),\quad T_* =\sup_xT(x),\quad
 T_0=T(0),\quad A=\tau^{*3}(0),\quad B=\tau^{*2}(0).
\end{equation}
All these quantities are finite at a fixed subcritical parameter.

\begin{lemma}\label{lem:budget}
Uniformly for $1\le p<p_c$ and sufficiently small $s>0$,
\begin{equation}\label{eq:budget}
 T_*=O(s),\qquad T_0=O(s^2),\qquad
 0\le A-1\le3T_0,\qquad0\le A-B\le3T_0,
 \qquad B=1+O(s^2).
\end{equation}
\end{lemma}

\begin{proof}
Set $h=\tau-\delta\ge0$. The first-edge connection bound gives
$h\le q*\tau$: for a nonzero endpoint choose the first bond of a simple
open path, remove that bond, and use independence and a union bound.
Since $\tau\ge\delta$,
\begin{equation}\label{eq:triple-positive}
 \tau^{*3}-\delta
 =h*(\tau^{*2}+\tau+\delta)
 \le3h*\tau^{*2}\le3q*\tau^{*3}=3T.
\end{equation}
Convolving once with $q$ shows
\[
 T(x)\le q(x)+3(q^{*2}*\tau^{*3})(x).
\]
The Fourier integrand $p^2d(k)^2\widehat\tau(k)^3$ is nonnegative.
Fourier inversion therefore bounds the last convolution at every $x$ by
its value at zero. By \eqref{eq:IR-simplified},
\begin{equation}\label{eq:U}
 U:=(q^{*2}*\tau^{*3})(0)
 \le\frac{b_s^3}{p}J_{2,3}=O(s^2).
\end{equation}
Thus $T_*\le p\|D\|_\infty+3U=O(s)$ by
\eqref{eq:coarse-pc}, while $T_0\le3U$ because $D(0)=0$.
Equation~\eqref{eq:triple-positive} at zero gives $A-1\le3T_0$.
Finally $1\le B\le A$, so both remaining assertions follow.
\end{proof}

In particular, the closed value $T_0$ is smaller than the spatial
supremum $T_*$. Keeping that distinction will turn the error terms below
into $O(s^3)$, whereas a bound by $T_*^2$ alone would give only $O(s^2)$.

\section{A uniform upper bound on the first lace coefficient}
\label{sec:pi1}

The decomposition below follows van der Hofstad and Sakai
\cite[Section~3.2, (3.45)--(3.55)]{HS2005}: nonzero first attachment,
terminal membership in a deleted-bond cluster, and deletion of the first bond.
The additional estimate is the separation $T_*=O(s)$ from
$T_0,A-1,A-B=O(s^2)$. It improves the remainder to $O(s^3)$ in
this changing-exponent regime. The reusable bound to keep in view is
\eqref{eq:exact-positive-bound}: an exact subtraction $p^2S_2$ and
nonnegative errors each carrying an extra small factor.

We specify the source events before estimating them. This also records
which restrictions survive the diagram bounds. Write $x\Leftrightarrow y$
if $x=y$ or there are two bond-disjoint open paths from $x$ to $y$.
For a set $\mathcal A$, the event
$\through{x}{y}{\mathcal A}$ means that $x$ and $y$ are connected and
every open path between them has a bond with an endpoint in $\mathcal A$,
or that $x=y\in\mathcal A$. This is the ``through'' convention of
\cite[Definition~3.1]{BCHSS2005}. Let $\widetilde{\mathcal C}^{\{u,v\}}(x)$
be the cluster of $x$ when bond $\{u,v\}$ is forced vacant.

A directed bond $(a,b)$ is pivotal for $v\leftrightarrow x$ if its
removal disconnects $v$ from $x$, with $a$ on the $v$ side and $b$ on
the $x$ side. Let $\mathcal P(v,x)$ be the set of these directed
pivotal bonds. Define
\begin{equation}\label{eq:Eprime}
 E'(v,x;\mathcal A)=\through{v}{x}{\mathcal A}
 \cap\left\{\nexists(a,b)\in\mathcal P(v,x):
                    \through{v}{a}{\mathcal A}\right\}.
\end{equation}
On the connection event in this definition the pivotal bonds are open.
The event $E'$ itself need not be increasing. We use the van den Berg--Kesten
(BK) inequality only on the
increasing connection events in its cited diagrammatic upper bounds.

Use subscripts $0,1,2$ for independent percolation configurations at the
same parameter. With the set from level~0 held fixed in the level~1
probability, the first coefficient is
\begin{equation}\label{eq:pi1-definition}
 \pi_p^{(1)}(x)=\sum_{u,v}q(v-u)\,
 \mathbb E_0\left[\ind_{\{0\Leftrightarrow u\}}
  \mathbb P_1\left(E'(v,x;
    \widetilde{\mathcal C}_0^{\{u,v\}}(0))\right)\right].
\end{equation}
This is \cite[(3.31)]{BCHSS2005}, also used in
\cite[Section~3.2, (3.44)--(3.45)]{HS2005}, with each bond weight
replaced by its probability $q(v-u)$. Set
$\overline\pi_p^{(1)}=\sum_x\pi_p^{(1)}(x)$.

\begin{proposition}\label{prop:pi1}
There is a constant $K_2$ such that, uniformly for
$1\le p<p_c$ and $0<s<\sigma_0$,
\begin{equation}\label{eq:pi1-bound}
 0\le\overline\pi_p^{(1)}
 \le(2p-p^2)S_2+K_2s^3\le S_2+K_2s^3.
\end{equation}
\end{proposition}

\subsection{The two restricted error terms}

For $u=0$ abbreviate
$\mathcal C_v=\widetilde{\mathcal C}_0^{\{0,v\}}(0)$ and
$\tau_v^{\del}(x)=\Pp(x\in\mathcal C_v)$. Split the sum in
\eqref{eq:pi1-definition} into the three disjoint cases
\begin{equation}\label{eq:split}
 u\ne0;\qquad u=0,\ x\notin\mathcal C_v;\qquad
 u=0,\ x\in\mathcal C_v.
\end{equation}
The last part is at most
\begin{equation}\label{eq:M}
 M_p=\sum_{v,x}q(v)\tau_v^{\del}(x)\tau(x-v),
\end{equation}
by dropping the no-pivotal requirement, retaining only the connection
$v\leftrightarrow x$ at level~1, and using independence of the two
levels.

To bound the first two parts, let $\circ$ denote disjoint occurrence
with disjoint open-bond witnesses. The source events
\cite[(4.2), (4.6)]{BCHSS2005} for $N=1$ are
\begin{align*}
 F_0(0,u,w,z)&=\{0\leftrightarrow u\}\circ\{0\leftrightarrow w\}
   \circ\{w\leftrightarrow u\}\circ\{w\leftrightarrow z\},\\
 F_1(v,t,z,x)&=\{v\leftrightarrow t\}\circ\{t\leftrightarrow z\}
   \circ\{t\leftrightarrow x\}\circ\{z\leftrightarrow x\}.
\end{align*}
Their terminal inclusion is
\begin{equation}\label{eq:terminal-inclusion}
 E'(v,x;\mathcal C)\subseteq
 \bigcup_{z\in\mathcal C}\ \bigcup_t F_1(v,t,z,x).
\end{equation}
In the joint event of \eqref{eq:pi1-definition}, the initial double
connection and the cluster membership in this inclusion are bounded by
$F_0$ at level~0, as in \cite[(4.10)]{BCHSS2005}. In particular the
vertex $u$ in $F_0(0,u,w,z)$ is the original $u$ of
\eqref{eq:pi1-definition}; it has not been summed out or renamed.
In words, the terminal inclusion retains a vertex in the specified cluster.
If the endpoint lies outside that cluster, these two vertices cannot coincide;
this strict exclusion must survive the later upper bounds.

Put $\widetilde\tau=q*\tau$ and
\[
 A_3(a,b,c)=\tau(c-a)\tau(b-a)\tau(c-b).
\]
The union bound, BK for $F_0,F_1$, and summation over $v$ give for the
$u\ne0$ part
\begin{equation}\label{eq:first-error-raw}
 \sum_{u\ne0,w} A_3(0,u,w)
 \sum_{z,t,x}\widetilde\tau(t-u)\tau(z-w)A_3(z,t,x).
\end{equation}
For each $u,w$, Cauchy--Schwarz gives
$(\tau*\tau)(r)\le B$, and therefore
\begin{align*}
 \sum_{z,t,x}\widetilde\tau(t-u)\tau(z-w)A_3(z,t,x)
 &=\sum_{z,t}\widetilde\tau(t-u)\tau(z-w)
     \tau(t-z)(\tau*\tau)(t-z)\\
 &\le B T(w-u)\le BT_*.
\end{align*}
Also $\sum_{u,w}A_3(0,u,w)=A$. The restriction $u\ne0$ excludes the
entire $u=0$ slice, whose sum is $\sum_w\tau(w)^2=B$. Consequently
\begin{equation}\label{eq:first-error}
 \text{the first part of \eqref{eq:split}}
 \le BT_*(A-B)=O(s^3).
\end{equation}

For the second part, \eqref{eq:terminal-inclusion} supplies a vertex
$z\in\mathcal C_v$. The condition $x\notin\mathcal C_v$ therefore
forces the precise restriction $z\ne x$. Retain it before dropping
$x\notin\mathcal C_v$. Since
$\Pp(z\in\mathcal C_v)\le\tau(z)$, BK bounds this part by
\begin{equation}\label{eq:second-error-raw}
 \sum_{v,z,t,x:\,z\ne x}q(v)\tau(z)\tau(t-v)
          \tau(t-z)\tau(t-x)\tau(z-x).
\end{equation}
For $r=t-z$, the strict terminal sum is
\begin{equation}\label{eq:hbad}
 H(r)=\tau(r)\big[(\tau*\tau)(r)-\tau(r)\big]\ge0,
 \qquad \sum_rH(r)=A-B.
\end{equation}
Indeed, the removed $x=z$ term has value $\tau(r)^2$.
Summing $v,z$ in \eqref{eq:second-error-raw} now gives
\begin{equation}\label{eq:second-error}
 \sum_r H(r)(q*\tau^{*2})(r)
 \le T_*(A-B)=O(s^3),
\end{equation}
where $q*\tau^{*2}\le q*\tau^{*3}$ follows from $\tau\ge\delta$.
Together with \eqref{eq:M}, this proves the uniform reduction
\begin{equation}\label{eq:pi1-M}
 \overline\pi_p^{(1)}\le M_p+O(s^3).
\end{equation}
All error constants in this reduction follow from
Lemma~\ref{lem:budget}; no estimate on a critical coefficient is used.

\subsection{The deleted bond and the main convolution}

Let $F_p=(q*\tau^{*2})(0)$. Deletion reduces connection probabilities,
so $M_p\le F_p$. A necessary improvement comes just from the terms
$x=v$. Conditioning on bond $\{0,v\}$ gives the exact identity
\begin{equation}\label{eq:bond-conditioning}
 \tau(v)=q(v)+[1-q(v)]\tau_v^{\del}(v)\qquad(v\ne0).
\end{equation}
The parameter range in use has $\|q\|_\infty=O(s)<1$. Retaining the
improvement at $x=v$, where $\tau(x-v)=1$, yields
\begin{align}
 M_p&\le F_p-\sum_vq(v)[\tau(v)-\tau_v^{\del}(v)]\notag\\
 &=F_p-p^2S_2+R_p,\qquad
 R_p=\sum_vq(v)^2\tau_v^{\del}(v).\label{eq:deleted-subtraction}
\end{align}
The nonnegative remainder is cubic:
\begin{equation}\label{eq:deleted-error}
 0\le R_p\le\|q\|_\infty\sum_vq(v)\tau(v)
 =\|q\|_\infty(q*\tau)(0)
 \le\|q\|_\infty T_0=O(s^3).
\end{equation}
This is an exact deleted-bond subtraction; expanding an inverse vacancy
probability is unnecessary.

Before making asymptotic estimates, the preceding event bounds give the
single explicit inequality
\begin{equation}\label{eq:exact-positive-bound}
 \boxed{\quad
 \overline\pi_p^{(1)}\le F_p-p^2S_2
  +\|q\|_\infty T_0+(B+1)T_*(A-B).
 \quad}
\end{equation}
The two terms after the subtraction are nonnegative and each is
$O(s^3)$ by Lemma~\ref{lem:budget}. Both strict restrictions leave
the factor $A-B$. More sharply, the positive first-edge inequality gives
$A-B=(h*\tau^{*2})(0)\le T_0$ and
$B-1\le2(q*\tau^{*2})(0)\le2T_0$.

To estimate $F_p$, use $h=\tau-\delta\le q*\tau$ and $D(0)=0$:
\begin{align}
 F_p&=2p(D*h)(0)+p(D*h*h)(0)\notag\\
 &\le2p^2(D^{*2}*\tau)(0)
      +p^3(D^{*3}*\tau^{*2})(0).\label{eq:F-decomposition}
\end{align}
The first Fourier integrand below is nonnegative. Applying
\eqref{eq:IR-simplified} with its factor $p$ retained gives
\begin{align}
 p^2(D^{*2}*\tau)(0)
 &\le p b_s I_{2,1},\label{eq:V2}\\
 0\le p^3(D^{*3}*\tau^{*2})(0)
 &\le p b_s^2 J_{3,2}.\label{eq:V3}
\end{align}
For \eqref{eq:V3} the convolution is nonnegative in real space, but
its Fourier integrand $p^3d(k)^3\widehat\tau(k)^2$ can have either sign.
Taking its absolute value before using \eqref{eq:IR-simplified} is what
justifies the bound.

Equations~\eqref{eq:I21}, \eqref{eq:needed-J}, and
\eqref{eq:coarse-pc} now imply
\[
 F_p\le2p b_s I_{2,1}+p b_s^2J_{3,2}
      =2pS_2+O(s^3).
\]
Combining this with \eqref{eq:pi1-M}--\eqref{eq:deleted-error} proves
\[
 \overline\pi_p^{(1)}\le(2p-p^2)S_2+O(s^3).
\]
Finally $2p-p^2=1-(p-1)^2\le1$. This proves
Proposition~\ref{prop:pi1}, including its uniformity over the full
subcritical interval $[1,p_c)$.

\section{The positive finite-order remainder and the threshold}
\label{sec:threshold}

We use the finite-order lace identity, rather than an infinite signed
series. Let
\[
 \pi_p^{(0)}(x)=\Pp(0\Leftrightarrow x)-\delta(x),\qquad
 \Pi_1=\pi_p^{(0)}-\pi_p^{(1)}.
\]
The identity at truncation order $M=1$ is
\begin{equation}\label{eq:finite-lace}
 \tau=\delta+q*\tau+\Pi_1*q*\tau+\Pi_1+\mathcal R_1,
 \qquad \mathcal R_1(x)\ge0.
\end{equation}
This is \cite[(3.9)--(3.10), (3.32)]{BCHSS2005}; the weighted
infinite-lattice version is \cite[Section~4.2, (4.6)]{HHS2008}.
For clarity, its positive remainder has the form
\begin{align}\label{eq:remainder-positive}
 \mathcal R_1(x)
 =\sum_{u,v,a,b}&q(v-u)q(b-a)\,
 \mathbb E_0\Big[\ind_{\{0\Leftrightarrow u\}}
 \mathbb E_1\Big[\ind_{E'(v,a;\widetilde{\mathcal C}_0^{\{u,v\}}(0))}
 \notag\\[-2pt]
 &\hspace{39mm}\times
 \mathbb P_2\big(\through{b}{x}{\widetilde{\mathcal C}_1^{\{a,b\}}(v)}\big)
 \Big]\Big].
\end{align}
Dropping the through-cluster constraint at level~2 gives the pointwise
bound from \cite[(3.36)]{BCHSS2005}:
\begin{equation}\label{eq:remainder-bound}
 0\le\mathcal R_1\le\pi_p^{(1)}*q*\tau.
\end{equation}

The cited construction applies to subcritical infinite graphs, where
clusters are almost surely finite, and its probability identities allow
inhomogeneous independent bond weights. We use HHS's weighted
infinite-volume identity directly. Connection-event BK bounds on the
countable bond set can equivalently be obtained by finite witness
approximation. No monotone convergence assertion for the nonmonotone
event $E'$ is needed. For each fixed $p<p_c$, all sums in
\eqref{eq:finite-lace} are absolutely convergent: $\pi_p^{(0)}\ge0$ and
\[
 \sum_x\pi_p^{(0)}(x)\le\chi(p)-1<\infty,
\]
Proposition~\ref{prop:pi1} bounds the sum of $\pi_p^{(1)}$, and
\eqref{eq:remainder-bound} gives
$\sum_x\mathcal R_1(x)\le p\overline\pi_p^{(1)}\chi(p)<\infty$.
The signed convolution is absolutely summable by these same bounds.

Set
\[
 U_s=S_2+K_2s^3,\qquad
 a_p=1+\sum_x\pi_p^{(0)}(x)-\overline\pi_p^{(1)}.
\]
Reduce $\sigma_0$ so that $U_s<1/2$. Nonnegativity of
$\pi_p^{(0)}$ and Proposition~\ref{prop:pi1} imply
$a_p\ge1-U_s>0$. Summing \eqref{eq:finite-lace} gives the exact scalar
identity
\begin{equation}\label{eq:scalar}
 \chi(p)[1-pa_p]=a_p+\sum_x\mathcal R_1(x)>0.
\end{equation}
Since $\chi(p)$ is positive and finite, for every $1\le p<p_c$,
\begin{equation}\label{eq:subcrit-upper}
 p<\frac1{a_p}\le\frac1{1-U_s}.
\end{equation}
If $p_c=1$ the same upper bound is automatic; otherwise take the supremum
over the scalar parameters $p\in[1,p_c)$. This yields
\begin{equation}\label{eq:pc-upper}
 p_c\le\frac1{1-S_2-K_2s^3}=1+S_2+O(s^3),
\end{equation}
because $S_2=O(s^2)$ and hence $U_s^2=O(s^4)$.
There is no passage to a critical Fourier transform, susceptibility
divergence, or a critical limit of a lace coefficient in this argument.
Avoidance of coefficient continuity is not a new general goal:
Liu \cite[Appendix~A]{Liu2025} gives a different argument at criticality
through a convolution identity and nonsummability. Here the positive finite
truncation gives a scalar inequality at every subcritical parameter. The saving
is in the additional correction calculation, not a replacement for the
lace-expansion work underpinning the imported infrared theorem.

The lower bound in Proposition~\ref{prop:NB} and
\eqref{eq:pc-upper} prove \eqref{eq:main}. Formula~\eqref{eq:S2}
then proves \eqref{eq:limit}. To make the uniform quantifiers explicit,
choose $\sigma_0$ smaller than $1/6$, the fixed smallness threshold needed
for \eqref{eq:HHS-assumptions}, and the fixed thresholds used for
$\|q\|_\infty<1$ and $U_s<1/2$. Every estimate above then holds for
every real $0<s<\sigma_0$ and every $p\in[1,p_c(s))$, with common
constants. This completes the proof of Theorem~\ref{thm:main}.

\section{A two-scale criterion and additional kernel families}
\label{sec:extensions}

The proof separates two scales rather than requiring a specific zeta kernel.
The following consequence makes that distinction reusable. All kernels in this
section define independent undirected bonds with the exact linear law $pD(x)$.

\begin{theorem}[Two-scale kernel criterion]\label{thm:criterion}
Let $D_\epsilon$ be even probability kernels on $\Z$, with $D_\epsilon(0)=0$,
positive nearest-neighbor weights, and $d_\epsilon(k)=\widehat D_\epsilon(k)<1$
for $0<|k|\le\pi$. Suppose the admissible percolating set is nonempty and,
with common constants as $\epsilon\downarrow0$,
\[
 \|D_\epsilon\|_\infty=O(\epsilon),\qquad
 \intk\frac{|d_\epsilon|^2}{(1-d_\epsilon)^3}\dm=O(\epsilon^2),
 \qquad
 \intk\frac{|d_\epsilon|^3}{(1-d_\epsilon)^2}\dm=O(\epsilon^3).
\]
Then
\[
 p_c(D_\epsilon)=1+\sum_xD_\epsilon(x)^2+O(\epsilon^3).
\]
The implicit constant is uniform over families with common hypothesis constants
and a common sufficiently small parameter interval.
\end{theorem}

\begin{proof}
Write $a=J_{2,3}$, $b=J_{3,2}$ and $m=\|D_\epsilon\|_\infty$.
Since $0\le1-d\le2$, $S_2=\int d^2\dm\le8a=O(\epsilon^2)$,
and the pointwise third moment is at most $mS_2=O(\epsilon^3)$.
The proof of Proposition~\ref{prop:NB} uses only evenness, summability
and these moment bounds. It supplies the lower bound with the stated error.
Its test parameters are admissible because $m\to0$.

Assumption~2.1 of \cite{HHS2008} holds with $\beta=O(\epsilon)$.
The same theorem used in Section~\ref{sec:diagrams} gives
$p_c=1+O(\epsilon)$ and \eqref{eq:IR-simplified} with
$b_\epsilon=1+O(\epsilon)$. Its dimension remains fixed at one.
The proof of Lemma~\ref{lem:budget} consequently gives
$T_*=O(\epsilon)$, $T_0=O(\epsilon^2)$ and
$A-1,A-B,B-1=O(\epsilon^2)$. Moreover,
\[
 |I_{2,1}-S_2|\le2J_{3,2}=O(\epsilon^3).
\]
The event bounds in Section~\ref{sec:pi1} use neither the power-law formula
nor positivity of its Fourier transform. With these new bounds their exact
deleted-bond inequality gives
$\overline\pi_p^{(1)}\le(2p-p^2)S_2+O(\epsilon^3)$ uniformly
subcritically. Section~\ref{sec:threshold} then gives
$p_c\le(1-S_2-O(\epsilon^3))^{-1}=1+S_2+O(\epsilon^3)$.
All imported results, and the inherited event decomposition, are the same as
for the original theorem; no new lace expansion is claimed.
\end{proof}

\begin{theorem}[Summable positive additions]\label{thm:additions}
Fix $B_0<\infty$. For any sequence $b_n\ge0$ with
$B=\sum_{n\ge1}b_n\le B_0$, define
\[
 D_s(\pm n)=\frac{n^{-1-s}+b_n}{2(Z_s+B)},\qquad n\ge1,\qquad D_s(0)=0.
\]
Uniformly over these sequences, for sufficiently small $s>0$,
\begin{align*}
 p_c&=1+\frac{\zeta(2+2s)+2\sum_{n\ge1}b_nn^{-1-s}
                    +\sum_{n\ge1}b_n^2}{2(Z_s+B)^2}+O_{B_0}(s^3)\\
 &=1+\frac{s^2}{2}\left(\zeta(2)+2\sum_{n\ge1}\frac{b_n}{n}
                                      +\sum_{n\ge1}b_n^2\right)
                                      +O_{B_0}(s^3).
\end{align*}
No moment assumption beyond summability is imposed on the additions.
\end{theorem}

\begin{proof}
Let $d_s^0$ denote the original power-law characteristic function.
The new kernel is a mixture with original-kernel weight
$w_s=Z_s/(Z_s+B)$. Thus
\[
 |d_s(k)|\le |d_s^0(k)|+B/Z_s
 \le s[L_0+B_0+\log_+(1/|k|)],
 \qquad 1-d_s(k)\ge w_s(1-d_s^0(k)).
\]
For small $s$ uniformly in $B\le B_0$, $w_s\ge1/2$, so the latter
is at least $\tfrac16(|k|/(2\pi))^s$. The integral proof of
Lemma~\ref{lem:fourier} applies with these modified constants, giving
$J_{r,q}=O_{B_0}(s^r)$ for fixed $r,q$ and $qs\le1/2$.
Also $\|D_s\|_\infty\le(1+B_0)s/2$. Full support off zero gives
the structural hypotheses. To check that admissible percolation exists,
at $p=3$ the model dominates the original kernel at parameter $3w_s>2$
for small $s$; Theorem~\ref{thm:main} places that model above its threshold.
The parameter $3$ is admissible since $3\|D_s\|_\infty<1$ eventually.
Theorem~\ref{thm:criterion} proves the first formula.

For the second, $\sum b_n^2\le B_0^2$ and
\[
 \left|\sum b_nn^{-1-s}-\sum b_n/n\right|
 \le s\sum b_n\frac{\log n}{n}=O(B_0s).
\]
Together with $\zeta(2+2s)=\zeta(2)+O(s)$ and
$(Z_s+B)^{-2}=s^2+O_{B_0}(s^3)$, this proves the assertion.
\end{proof}

For example, $b_1=\lambda\ge0$ and $b_n=0$ for $n>1$ changes the
coefficient from $\pi^2/12$ to $\pi^2/12+\lambda+\lambda^2/2$.
Thus the coefficient is sensitive to local connection weights even though the
long-distance exponent is unchanged.

\begin{theorem}[Compact mixtures of exponents]\label{thm:mixtures}
Fix $0<c_-\le c_+<\infty$. For any probability measure $\nu$ on
$[c_-,c_+]$, let $D_s=\int D^0_{cs}\,\nu(\mathrm dc)$, where
$D^0_t$ is the power-law kernel~\eqref{eq:kernel}. Uniformly over $\nu$,
\[
 p_c(D_s)=1+\frac{\pi^2}{12}\left(\int c\,\nu(\mathrm dc)\right)^2s^2
           +O_{c_-,c_+}(s^3).
\]
\end{theorem}

\begin{proof}
For $c_+s\le1/2$, integration of Lemma~\ref{lem:fourier} gives
\[
 |d_s(k)|\le c_+s[L_0+\log_+(1/|k|)],\qquad
 1-d_s(k)\ge\tfrac13(|k|/(2\pi))^{c_+s}.
\]
Thus $J_{r,q}=O(s^r)$ when $qc_+s\le1/2$, and
$\|D_s\|_\infty\le c_+s/2$. Full support again gives the structural
conditions. For admissible percolation, integral comparison yields
\[
 \frac{D^0_{cs}(n)}{D^0_{c_+s}(n)}
 =\frac{Z_{c_+s}}{Z_{cs}}n^{(c_+-c)s}
 \ge\frac{c_-}{c_+(1+c_+s)}\ge\frac{c_-}{2c_+}
\]
for small $s$. The fixed parameter $p=6c_+/c_-$ is eventually admissible
and dominates an original-kernel model at parameter at least $3$.
Apply Theorem~\ref{thm:main}, then Theorem~\ref{thm:criterion}.
Finally Tonelli's theorem gives
\[
 S_2=\iint\frac{\zeta(2+(c+d)s)}{2Z_{cs}Z_{ds}}
                  \,\nu(\mathrm dc)\nu(\mathrm dd).
\]
Uniformly on the compact parameter square, $Z_{cs}^{-1}=cs+O(s^2)$
and $\zeta(2+(c+d)s)=\zeta(2)+O(s)$. This proves the formula.
\end{proof}

\section{Scope and further questions}

The correction in Theorem~\ref{thm:main} is the two-step return weight
of the exact normalized kernel. The proof requires both the uniform
small-exponent Fourier bounds and the percolation event restrictions;
evaluating random-walk return probabilities alone would not identify
the lace coefficient. The positive odd-truncation remainder is enough
for an upper bound on the threshold, so a matching asymptotic formula
for $\overline\pi_p^{(1)}$ and a bound on every higher coefficient are
not part of the argument.

The result fixes dimension one and lets only $\sigma$ tend to zero.
The error constant is not evaluated numerically, and the proof gives
neither a third-order coefficient nor a useful quantitative estimate
throughout $0<\sigma<1$. Other bond laws, including exponential
connection probabilities, may differ at the correction scale and are
not covered by the theorem. The argument also does not determine the
oriented critical correction. The uniform kernel estimates and the
one-sided lace estimate yield the sufficient criterion and the two additional
families in Section~\ref{sec:extensions}. Signed additions, unbounded exponent
ratios and useful numerical error constants remain outside those statements.

\section*{Declarations}

\paragraph{Data and reproducibility.}
This is an analytic study. No numerical experiment or empirical data
is used in the proof. The model, estimates, and imported results are
specified in the text; the manuscript source and bibliography form
the associated research materials.

\paragraph{Ethics.}
The work involves no human participants, animals, personal data, or
experimental intervention.

\paragraph{Contributions and AI disclosure.}
OpenAI Codex was used to develop the arguments, compare the cited
sources, perform internal mathematical checks, and prepare this draft.
The scholarly attribution is Anonymous. Evidence Press is the publisher;
account ownership is not a claim of human mathematical contribution.
The supplied review and current checks are AI-assisted scrutiny, not external
human peer review or formal verification.

\paragraph{Funding and competing interests.}
No external funding or competing interests have been reported in the
materials supplied for this work. No institutional affiliation or funding
award is claimed. This reports the available record, not independent
verification of an anonymous creator's circumstances.

\paragraph{Archive boundary.}
The original submission contains internal file-hash bindings but omits four
historically cached external-source files and its original packaging test suite.
Those historical bytes and executions cannot be reconstructed from that ZIP.
Current source checking and any new diagnostics are recorded separately;
neither is a replacement for the missing historical records.

\begin{thebibliography}{9}

\bibitem{Amit2026}
G.~Amit,
\emph{Explicit upper bounds on the threshold of one-dimensional long-range
percolation},
arXiv:2610.11872v1 [cond-mat.stat-mech], 2026.
\url{https://arxiv.org/abs/2610.11872v1}.

\bibitem{BCHSS2005}
C.~Borgs, J.~T.~Chayes, R.~van der Hofstad, G.~Slade, and J.~Spencer,
\emph{Random subgraphs of finite graphs. II. The lace expansion and the
triangle condition},
Annals of Probability \textbf{33} (2005), no.~5, 1886--1944.
\href{https://doi.org/10.1214/009117905000000260}{doi:10.1214/009117905000000260}.
arXiv:math/0401070.

\bibitem{HHS2008}
M.~Heydenreich, R.~van der Hofstad, and A.~Sakai,
\emph{Mean-field behavior for long- and finite range Ising model, percolation
and self-avoiding walk},
Journal of Statistical Physics \textbf{132} (2008), no.~6, 1001--1049.
\href{https://doi.org/10.1007/s10955-008-9580-5}{doi:10.1007/s10955-008-9580-5}.
arXiv:0712.0312v2.

\bibitem{HS2005}
R.~van der Hofstad and A.~Sakai,
\emph{Critical points for spread-out self-avoiding walk, percolation and the
contact process above the upper critical dimensions},
Probability Theory and Related Fields \textbf{132} (2005), no.~3, 438--470.
\href{https://doi.org/10.1007/s00440-004-0405-4}{doi:10.1007/s00440-004-0405-4}.
arXiv:math/0402050v2.

\bibitem{Liu2025}
Y.~Liu,
\emph{High-dimensional long-range statistical mechanical models have random
walk correlation functions}, arXiv:2502.12104v3, 2025.
\url{https://arxiv.org/abs/2502.12104v3}.

\end{thebibliography}
\end{document}
